\chapter{Introduction}
\label{chap:introduction}

\section{Background and Motivation}

Large language models (LLMs) are now widely used by university students to draft literature reviews, summarize papers, and format reference lists. Although these tools improve writing speed, they also introduce a severe reliability problem known as \emph{citation hallucination}~\cite{Ji2023}. Because generative models predict probable word sequences rather than querying external databases, they frequently output references that look syntactically valid yet fail basic scholarly verification. Walters and Wilder~\cite{walters2023fabrication} showed that a substantial share of citations generated by conversational LLMs either do not exist or contain fabricated Digital Object Identifiers (DOIs), wrong publication years, or altered author lists.

Manual verification of references does not scale in undergraduate and graduate education. A typical graduation project report cites between thirty and eighty sources across journals, conference proceedings, books, web documentation, and national legal statutes. Checking each reference by hand requires searching multiple databases, comparing author names and publication years, verifying whether a DOI resolves to the claimed title, and reading the source abstract to confirm that the cited claim matches the paper. Furthermore, citation errors take several distinct forms that carry very different academic meanings:
\begin{enumerate}
    \item \textbf{Fabricated or chimeric sources:} The cited paper does not exist in any scholarly index, or a real DOI from a famous article is attached to an invented title.
    \item \textbf{Metadata or claim errors:} The underlying work is real, but the student recorded the wrong year, omitted co-authors, cited a retracted paper, or attached a body claim that contradicts the source abstract.
    \item \textbf{Structural linking errors:} An in-text citation marker such as \texttt{(Vaswani et al., 2017)} has no matching bibliography entry, or bibliography entries are listed to inflate the reference count without ever being cited in the text.
\end{enumerate}

Two engineering principles guide the solution developed in this graduation project. First, no single bibliographic registry indexes the entire scholarly record~\cite{visser2021bibliographic}. Computer science workshop papers often appear in DBLP or the ACL Anthology before Crossref, biomedical articles live in PubMed and Europe PMC, and Vietnamese academic reports frequently cite national laws (\begin{vietnamese}\textit{Luật}, \textit{Nghị định}\end{vietnamese}). Reliable verification therefore requires federated retrieval across multiple complementary databases. Second, an automated tool used in academic evaluation must never turn network failures or ambiguous matches into false accusations. Following the principle of selective classification~\cite{chow1970optimum}, \cite{Geifman2017}, the system must explicitly abstain when external APIs fail or when similarity scores fall into an uncertain border zone, leaving final judgment to the human instructor.

\section{Problem Statement and Title Scope}

The registered title of this graduation project is \textit{``Automated Essay Grading and Source Hallucination Detection via Neuro-Symbolic Logic Checker''}. In educational software engineering, holistic essay grading includes subjective qualities such as argument structure, rhetorical style, and grammatical fluency. Grading those subjective traits with black-box neural models lacks transparency and is easily disputed by students. Therefore, this project deliberately scopes the \emph{automated grading} component to an objective, verifiable sub-dimension of academic assessment: \textbf{Automated Citation Integrity Scoring (CIS)}.

Formally, given an academic manuscript submitted in Portable Document Format (\texttt{.pdf}), Office Open XML (\texttt{.docx}), BibTeX (\texttt{.bib}), or Research Information Systems (\texttt{.ris}) format, the software system must solve four connected tasks:
\begin{enumerate}
    \item \textbf{Multi-format extraction and bidirectional linking:} Extract in-text citation markers (with surrounding sentence context) and bibliography entries across diverse formatting styles, then reconstruct the bipartite graph linking body markers to reference entries.
    \item \textbf{Federated source and claim verification:} Query offline and live scholarly registries, verify identifier liveness and retraction status, align in-text claims against source abstracts across English and Vietnamese, and classify each reference into a clear five-class taxonomy.
    \item \textbf{Calibrated scoring and bibliography repair:} Aggregate structural and bibliographic evidence into an interpretable $[0, 100]$ Citation Integrity Score (CIS) with calibrated confidence probabilities, and generate corrected reference entries from authoritative metadata.
    \item \textbf{Auditable human-in-the-loop review:} Present every verdict alongside its matched registry record and triggered logic rules inside a role-based web workspace where authorized teachers can inspect, override, and audit decisions.
\end{enumerate}

\section{Research Objectives}
\label{sec:objectives}

To address the problem statement above, this project defines six concrete engineering and empirical objectives:
\begin{itemize}
    \item \textbf{OBJ-1 (Multi-Format Extraction and Bidirectional Linking):} Build a hybrid document ingestion pipeline combining GROBID \texttt{0.9.1-crf} TEI-XML, two-column PyMuPDF, \texttt{python-docx}, and twelve modular citation-style parsers with a seven-method bidirectional citation linker.
    \item \textbf{OBJ-2 (Domain-Aware Federated Retrieval):} Construct an offline SQLite FTS5 scholarly index paired with an asynchronous orchestrator across more than fifteen public scholarly and technical APIs, Retraction Watch screening, and SSRF-safe URL validation.
    \item \textbf{OBJ-3 (Explainable Neuro-Symbolic Verification and Calibration):} Combine Sentence-BERT and ONNX INT8 Vietnamese-to-English semantic encoders with a priority-ordered symbolic rule engine, Multinomial Logistic Regression, and logit-space Platt probability calibration.
    \item \textbf{OBJ-4 (Interpretable Integrity Grading and Automated Remediation):} Formulate the five-component Citation Integrity Score (CIS), verify its rank stability under weight perturbations, and support automated bibliography repair into standard formats.
    \item \textbf{OBJ-5 (Secure Role-Based Web Platform):} Engineer a containerized FastAPI and React single-page application (\textbf{CITRA}) supporting Student, Teacher, and Administrator workflows, real-time Server-Sent Events (SSE) progress streaming, and defense-in-depth authorization and SSRF security controls.
    \item \textbf{OBJ-6 (Leak-Free Large-Scale Empirical Evaluation):} Evaluate the platform against six baselines (\texttt{B0}--\texttt{B5}) on a paper-disjoint 120,000-sample benchmark, a 2,400-sample out-of-distribution (OOD) stress suite, real multi-style PDF corpora, and 50-user concurrency load tests.
\end{itemize}

\section{System Scope and Operational Boundaries}

\subsection{In-Scope Capabilities}

Version \texttt{3.90.5} of \textbf{CITRA} supports English and Vietnamese academic manuscripts in \texttt{.pdf}, \texttt{.docx}, \texttt{.bib}, and \texttt{.ris} formats. It parses eleven international citation styles (APA, IEEE, ACM, Harvard, Vancouver, Nature, CSE, ACS, AMA, Chicago, and MLA) plus Vietnamese statutory documents (\begin{vietnamese}\texttt{Luật}, \texttt{Nghị định}, \texttt{Thông tư}, \texttt{Quyết định}\end{vietnamese}). Source verification covers journal articles, conference papers, preprints, books, open-access repositories, technical standards (IETF RFCs, Python packages, Hugging Face artifacts), and web resources. All services run on standard CPU hardware via Docker Compose without requiring a dedicated GPU or paid commercial database subscriptions.

\subsection{Out-of-Scope Boundaries}

First, \textbf{CITRA} does not grade prose quality, grammar, or topical originality, nor does it perform full-text copy-paste plagiarism detection against closed student essay repositories. Second, when a cited commercial article is behind a publisher paywall that hides both its full text and its abstract, claim-level semantic verification is restricted to available metadata and open abstracts rather than bypassing access controls. Third, system outputs represent automated decision-support evidence rather than automated disciplinary rulings; instructors retain full authority to override any machine verdict.

\section{Main Contributions}

Aligned with Objectives \textbf{OBJ-1} through \textbf{OBJ-6}, this graduation project delivers six technical contributions:
\begin{enumerate}
    \item \textbf{Hybrid Neuro-Symbolic Verification Engine:} Integrates Sentence-BERT (\texttt{all-MiniLM-L6-v2}) and ONNX INT8 multilingual embeddings with Order 0 resource pre-screening, 16 priority-ordered symbolic rule groups (68 distinct \texttt{R-*} rule identifiers), a 16-feature Multinomial Logistic Regression classifier ($C = 2.0$), and logit-space Platt scaling.
    \item \textbf{Provenance-Aware Selective Abstention:} Separates confirmed registry refutation (\texttt{suspected\_hallucination}, requiring clean HTTP 200 empty responses or exposed chimeric/fabricated identifiers) from network exhaustion (\texttt{HTTP 429}), regional DOIs lacking public metadata (\texttt{R-DOI-RA-NO-CSL}), and borderline similarity ($[0.40, 0.60)$), routing uncertain cases safely to \texttt{unresolved}.
    \item \textbf{Twelve-Style Extraction and Seven-Method Bidirectional Linker:} Combines GROBID \texttt{0.9.1-crf} TEI-XML pointers with two-column PyMuPDF fallback, table-of-contents filtering, Vietnamese statutory verification (\texttt{VNLegalVerifier}), and a bipartite linker detecting eight structural mapping states.
    \item \textbf{Tiered Offline-First Retrieval Architecture:} Pairs a 129,125-record local SQLite FTS5 corpus (128,105 unique papers; 82,877 abstracts) and 101 seminal papers with rate-limited, circuit-broken routing across 15+ live APIs and a pinned-IP SSRF-safe HTTP/Wayback validator.
    \item \textbf{Production-Grade Full-Stack Architecture (\textbf{CITRA}):} Implements an asynchronous FastAPI backend (15 SQLAlchemy ORM tables, 13 SHA-256 checksummed migrations (\texttt{0001}--\texttt{0013}), \texttt{PipelineAdmissionController}, and transactional CIS recalculation) and a responsive React 18.3 / TypeScript 5.8 web workspace verified by 2,320+ backend tests, 487+ Vitest tests (63 suites), and 20 Playwright E2E suites.
    \item \textbf{Reproducible Paper-Disjoint Benchmark Suite:} Establishes a zero-overlap 120,000-sample benchmark and a 2,400-sample OOD guardrail benchmark with paired McNemar, Holm--Bonferroni, and bootstrap significance protocols, achieving $99.02\%$ in-distribution accuracy ($\text{ECE} = 0.0395$) and $99.75\%$ OOD accuracy.
\end{enumerate}

\section{Report Organization}

The remainder of this graduation project report is organized into five chapters and five appendices, as summarized in Table~\ref{tab:outline}.

\begin{longtable}{|L{0.24}|L{0.76}|}
\caption{Organization of the graduation project report by chapter and appendix.}\label{tab:outline}\\
\hline
\textbf{Part} & \textbf{Summary of Contents} \\
\hline
\endfirsthead
\multicolumn{2}{l}{\itshape Table \thetable\ (continued from previous page)}\\
\hline
\textbf{Part} & \textbf{Summary of Contents} \\
\hline
\endhead
\hline
\multicolumn{2}{r}{\itshape Continued on next page}\\
\endfoot
\hline
\endlastfoot
Chapter 1 & Background, motivation, CIS grading scope, research objectives, system boundaries, and main contributions. \\
\hline
Chapter 2 & Literature review, comparison with related tools, scholarly registries, and mathematical foundations of semantic matching and calibration. \\
\hline
Chapter 3 & IEEE 29148 functional and ISO/IEC 25010 non-functional requirements, UML use cases, sequence workflow, 15-table data model, API surface, and security design. \\
\hline
Chapter 4 & Implementation of the eight software layers, design patterns, 68-rule symbolic engine, 16D ML + Platt calibration, CIS rubric, and React web interface. \\
\hline
Chapter 5 & Requirements traceability, automated test suites, experimental results across five research questions (\textbf{RQ1}--\textbf{RQ5}), and threats to validity. \\
\hline
Chapter 6 & Achievement mapping against objectives, decision-path summary, current system limitations, and future engineering directions. \\
\hline
Appendices A--E & Sanitized environment and Docker configurations, five-class taxonomy rules, package map, and project execution and contribution log. \\
\end{longtable}
